\section{Conclusion}

As part of this study, the research question was whether publicly available data can give enough information about the demand of features to assist product team decision-making. Once the framework had been worked out, the forecasts had been carried out and compared with the actual trends, the evidence points towards the affirmative conclusion, with certain proper caveats.

With the empirical case study of Notion, this research showed that the signals of search interest, review sentiment and patterns of feature mention hold valuable information on demand patterns. Coherent patterns were discovered when applying forecasting models to such features as Notion AI, templates, and workflow automation. The exponential growth pattern observed with the AI capabilities is in line with the available technological trends. Templates had uniform and trusted demand---a basic feature with predictable reliance on the user. The interest in automation rose at a slower rate, following larger trends in the productivity software industry.

These predictions have a practical value because they are coupled to design decisions. It is not useful predicting that the interest in Notion AI will increase by 30\% next quarter. How the action product teams should react to that information is the substantive question. That question is answered in this framework because forecast trajectories are converted into recommendations---prioritize this feature, maintain that one, reconsider this other. It acts as an intermediary between product management activities and the output of data science.

It is appropriate to give honest recognition to limitations. This framework fails to cover every situation. Any forecast can be nullified due to unexpected changes in the market, unanticipated competitive action, or outbreaks. External data is a proxy and not a measure of how active the internal users base is. The models that are used though valid have no predictive magic and this means that they will work best when historical patterns are there which is not always the case.

However, the bigger picture is that there is. The use of proprietary analytics infrastructure is not a requirement in the incorporation of data in feature prioritization processes. Google Trends can be used by smaller teams and startups, as well as individual researchers, to analyze app reviews and create forecasting models. The tools exist. What was lacking before was a framework that shows how to combine these elements when it comes to decision making of SaaS products. Such is the contribution that this work makes.

In future, the use of demand forecasting will probably become a more usual activity in software development. With the market of the SaaS becoming increasingly more competitive, and the user demands ever-growing, the foresight of understanding the user needs, as opposed to the response to the existing feedback, is a competitive edge. This paper is an attempt to offer product teams a platform on which to build that capacity. The key point to note is that it is important to study the future demand just as much as it is important to study the present behaviour. Through forecasting in the product development cycle, the organisation is able to create features that predict user needs instead of gearing features to it. This is a better way of developing software.
